\section{Security of Complete Stopped Texts}
\label{sec:stopping}

\subsection{A Public Stopping Likelihood}

RRC's published Algorithm 3 rotates its current point by an independent uniform
offset, selects a token, and stops when the point lies within one half-unit of
the selected interval's midpoint~\cite{yan2026rrc}. Conditional on the active
history and selected token, the rotated point is uniform in that interval.
For $B$ input bits and observed tokens $y_{1:t}$, define
\begin{equation}
 W_i=2^B\prod_{j=1}^{i}q_j(y_j\mid y_{<j}),\qquad
 h_i=\min\{1,W_i^{-1}\}.
\end{equation}
Here $q_i$ is the specified token source, which may already include processing.
The midpoint window has length one, so $h_i$ is the conditional probability of
stopping at step $i$. Fresh independent rotations then give the stopped law
\begin{equation}
 S(y_{1:t})=
 \left(\prod_{i=1}^{t}q_i(y_i\mid y_{<i})\right)
 \left(\prod_{i<t}(1-h_i)\right)h_t.
 \label{eq:stopped-law}
\end{equation}
The extra factor is computable from the public model, prompt, output, and $B$.
It does not require the rotated point or key.

Suppose a natural control first draws a length $L$ independently from
$\ell(t)=\Pr_S(\tau=t)$ and then generates $L$ tokens from $q$. Its law is
$C(y_{1:t})=\ell(t)\prod_i q_i(y_i\mid y_{<i})$. The length histograms agree
exactly, yet the likelihood ratio still contains
$\prod_{i<t}(1-h_i)h_t/\ell(t)$. This is content--length dependence, not a
length-only attack. Using the \emph{same content-dependent stopping law} for
the control would instead reproduce the ideal process; the comparator must
therefore be explicit.

\subsection{An Exact Counterexample}

Let $B=1$, $q(A)=3/4$, and $q(B)=1/4$. Table~\ref{tab:rrc-exact} lists the entire
stopped law, calculated with rational arithmetic. Its length probabilities are
$3/4$, $11/48$, and $1/48$ for lengths one, two, and three. An independently
length-matched natural source assigns probability $3/16$ to the complete text
$B$, whereas the stopped encoder assigns $1/4$. That event alone distinguishes
by $1/16$ with ideal randomness. Enumerating all natural strings at these three
lengths gives total variation $33/256$.

\begin{table}[t]
\centering
\caption{Exact stopped law of the published ideal RRC rule. The independent
control has the same marginal length distribution.}
\label{tab:rrc-exact}
\begin{tabular}{lrr}
\toprule
Complete text & Stopped law & Independent control \\
\midrule
$A$ & $1/2$ & $9/16$ \\
$B$ & $1/4$ & $3/16$ \\
$AA$ & $1/6$ & $33/256$ \\
$AB$ & $1/16$ & $11/256$ \\
$AAA$ & $1/64$ & $9/1024$ \\
$AAB$ & $1/192$ & $3/1024$ \\
Other strings & $0$ & positive \\
\bottomrule
\end{tabular}
\end{table}

This counterexample addresses the step from active next-token preservation to
complete-text indistinguishability in the paper's Appendix B.3. It does not
challenge the fresh-rotation sampling argument. Current RRC code adds reverse
verification before stopping, which changes the stopping law; neither the
formula nor the exact table is attributed to that revised implementation.

\subsection{The Gap Can Grow with Message Length}

The issue is not confined to a one-bit example. Consider a fixed binary IID
source with probabilities $p$ and $1-p$, where $0<p<1$ and $p\ne1/2$. Write
$S_n=-\sum_{i=1}^n\log_2 q(Y_i)$ for cumulative surprisal. The ideal RRC stopping
rule accepts at step $n$ with probability $\min(1,2^{S_n-B})$.

\begin{theorem}[Stopping separation]
\label{thm:stopping}
Let $P_B$ be the stopped-transcript law under this rule, and let $C_B$ draw the
exact marginal stopping length independently before sampling fresh IID tokens.
For every fixed nonuniform binary source,
\begin{equation}
 \operatorname{TV}(P_B,C_B)\ge
 1-O_p\!\left(\frac{\log B}{\sqrt B}\right)\longrightarrow1.
\end{equation}
\end{theorem}

\paragraph{Proof sketch}
Let $a$ and $b$ be the minimum and maximum symbol surprisals. Stopping is certain
once $S_n\ge B$, hence $S_\tau\le B+b$. Before $S_n<B-k$, the possible hazards
sum to at most $2^{-k}/(1-2^{-a})$, because the surprisal increases by at least
$a$ each step. Thus a stopped text lies in the band $[B-k,B+b]$ with high
probability. For an independently sampled length, the surprisal is a shifted
and scaled binomial variable. Its probability of occupying a band of width
$k+b$ is $O_p((k+1)/\sqrt B)$, except for exponentially unlikely very short
lengths. Choosing $k=\log_2 B$ proves the claim. Appendix~\ref{app:proofs}
gives constants and the elementary anticoncentration argument.

The assumption $p\ne1/2$ matters. For a uniform binary source, interval width
and hazard depend only on the step index. Stopping length is then independent
of token identities, and the independently length-matched law is identical.
The theorem concerns a fixed source as $B$ grows, not a source approaching
uniformity with $B$, and it does not assert monotonicity at every finite $B$.
Our use of information random walks follows established source-coding
analysis~\cite{drmota2010tunstall}; the security question here is the separation
between complete transcript laws after length matching.
